% ═══════════════════════════════════════════════════════════════
% PLANTILLA — Texto corto (Editorial, Carta al Editor)
%
% Contenido: el editorial de Aldana Lopez, Vol. 5, No. 2, 2025. Un
% articulo nuevo parte de este archivo, nunca de cero (LEEME.md). El
% texto de los autores no se toca; los comentarios son del taller.
%
% Usa: apm-editorial.cls (debe estar en el mismo directorio)
% Compila: bash taller/componer.sh taller/ejemplo_editorial.tex
% ═══════════════════════════════════════════════════════════════
\documentclass{apm-editorial}

% ── METADATA ─────────────────────────────────────────────────
\APMtitleEN[Neuromodulación en psiquiatría: un llamado a la renovación profesional e institucional]{Neuromodulaci\'on en la psiquiatr\'ia\\Un llamado a la renovaci\'on profesional e institucional}
\APMauthor{Jes\'us Alejandro Aldana L\'opez}
\APMaffiliation{Cl\'inica de Psiquiatr\'ia, Departamento de Cl\'inicas,
  Centro Universitario de Tlajomulco,\par
  Universidad de Guadalajara, Jalisco, M\'exico.}
\APMorcid{https://orcid.org/0009-0004-1234-5678}
\APMemail{jesus.aldana@academicos.udg.mx}
\APMphone{+52 33 1234 5678}
\APMaddress{Av. Ejido 200, Col. El Baj\'io, Tlajomulco de Z\'u\~niga, Jalisco, M\'exico, C.P. 45645}
\APMdates{11 marzo 2025}{27 abril 2025}{30 abril 2025}
\APMdoi{https://doi.org/10.5281/zenodo.17567781}
\APMauthorshort{Aldana L\'opez}
\APMauthorpdf{Jesús Alejandro Aldana López}
\APMkeywordspdf{Neuromodulación, Psiquiatría, Estimulación Magnética Transcraneal, Estimulación Cerebral Profunda, América Latina}
\APMtitleshort{Neuromodulaci\'on en psiquiatr\'ia}
% \APMfolio{...}  <- COMANDO MUERTO: el setter existe en el .cls pero el valor
%                    nunca se renderiza (el Folio fue eliminado del diseno).
%                    Ver norma/09_limitaciones_conocidas.md, 3a. No lo uses.
\APMvolume{5}{2}{2025}
\APMperiod{MAYO--AGOSTO}
\APMlogo{logo_hires.png}
% \APMlogoSixty{logo_60anos.png}  <- COMANDO MUERTO, mismo caso: el logo del
%                    60 aniversario no tiene codigo de render (norma/09, 3b).
%                    Comentado el 1 de septiembre de 2026; el PDF no cambia.

% ── ABSTRACT ─────────────────────────────────────────────────
\APMabstract{Este editorial hace un llamado a la integraci\'on inmediata de la neuromodulaci\'on en la educaci\'on y la pr\'actica psiqui\'atrica. Insta a los profesionales e instituciones de M\'exico y Am\'erica Latina a actualizar sus est\'andares, abordar las lagunas \'eticas y regulatorias, y asegurar una implementaci\'on responsable y generalizada que mantenga la integridad moral y cient\'ifica de la psiquiatr\'ia. Los formuladores de pol\'iticas deben desarrollar gu\'ias regulatorias estandarizadas, fomentar alianzas institucionales con organizaciones globales de neuromodulaci\'on y exigir programas de capacitaci\'on certificada para preparar a los cl\'inicos ante los avances emergentes. El establecimiento de redes de neuromodulaci\'on entre hospitales y cl\'inicas facilitar\'a el intercambio r\'apido de experiencias pr\'acticas y apoyar\'a la resoluci\'on colaborativa de problemas, acelerando as\'i la adopci\'on y mejorando los resultados cl\'inicos.}
\APMkeywords{Neuromodulaci\'on, Psiquiatr\'ia, Estimulaci\'on Magn\'etica
  Transcraneal, Estimulaci\'on Cerebral Profunda, Am\'erica Latina}

% ── COLOPHON ─────────────────────────────────────────────────
\APMconflict{El autor declara no tener conflicto de intereses.}
\APMfunding{Esta investigaci\'on no recibi\'o financiamiento
  espec\'ifico de agencias del sector p\'ublico, comercial o sin
  fines de lucro.}
\APMlicense{\textcopyright\ 2025 El Autor. Publicado bajo Creative
  Commons Atribuci\'on-No Comercial 4.0 Internacional (CC~BY-NC~4.0).}

% ═════════════════════════════════════════════════════════════
\begin{document}
\makeAPMeditorial
\pdfinfo{
  /Author (Jes\string\372s Alejandro Aldana L\string\363pez)
  /Keywords (Neuromodulaci\string\363n, Psiquiatr\string\355a, Estimulaci\string\363n Magn\string\351tica Transcraneal, Estimulaci\string\363n Cerebral Profunda, Am\string\351rica Latina)
}

% ── BODY ─────────────────────────────────────────────────────

{\noindent\fontsize{9.5}{12}\selectfont\bfseries\color{burg}INTRODUCCI\'ON}\par
\vspace{4pt}\dropcap{L}{a}
psiquiatr\'ia se encuentra en un punto de inflexi\'on hist\'orico.
La convergencia de la neurociencia, la tecnolog\'ia y la pr\'actica
cl\'inica ha dado lugar a un arsenal terap\'eutico cada vez m\'as
sofisticado. Entre las innovaciones m\'as prometedoras se encuentra la
neuromodulaci\'on, un campo que abarca un conjunto de t\'ecnicas
dise\~nadas para modificar la actividad neuronal de manera selectiva y
reversible. Desde la Estimulaci\'on Magn\'etica Transcraneal (EMT)
hasta la Estimulaci\'on Cerebral Profunda (ECP), estas herramientas
est\'an redefiniendo los l\'imites de lo posible en el tratamiento de
trastornos mentales graves y resistentes.

Sin embargo, el entusiasmo por el potencial de la neuromodulaci\'on
debe ir acompa\~nado de una reflexi\'on profunda y una acci\'on
decidida. La integraci\'on exitosa de estas terapias en la pr\'actica
cl\'inica diaria no es simplemente una cuesti\'on de adquirir nuevos
equipos, sino que exige una renovaci\'on fundamental a nivel
profesional e institucional. Este editorial es un llamado a la
comunidad psiqui\'atrica para abordar proactivamente los desaf\'ios que
esta nueva era nos presenta.

\section*{La Brecha Formativa: De la Teor\'ia a la Pr\'actica}

El primer desaf\'io es la formaci\'on. El psiquiatra contempor\'aneo
debe ser m\'as que un experto en psicofarmacolog\'ia y psicoterapia;
debe poseer una comprensi\'on s\'olida de la neuroanatom\'ia funcional,
los principios de la electrofisiolog\'ia y la f\'isica detr\'as de las
t\'ecnicas de neuromodulaci\'on. Los programas de residencia y los
cursos de educaci\'on m\'edica continua deben actualizarse urgentemente
para incluir m\'odulos te\'oricos y pr\'acticos rigurosos sobre estas
modalidades. No podemos permitir que la aplicaci\'on de estas potentes
herramientas quede en manos de personal insuficientemente capacitado o,
peor a\'un, que se convierta en un nicho comercial desregulado. La
certificaci\'on y la acreditaci\'on deben ser est\'andares, no opciones.

\section*{Desaf\'ios \'Eticos y Acceso Equitativo}

En segundo lugar, debemos enfrentar los dilemas \'eticos. ?`Qui\'enes
son los candidatos ideales para estas terapias? ?`C\'omo garantizamos
un consentimiento verdaderamente informado cuando los mecanismos de
acci\'on a\'un no se comprenden por completo? Y, quiz\'as lo m\'as
importante, ?`c\'omo aseguramos un acceso equitativo? El alto costo y
la disponibilidad limitada de la neuromodulaci\'on corren el riesgo de
ampliar la brecha de atenci\'on en salud mental, creando un sistema de
dos niveles donde solo los privilegiados pueden acceder a los
tratamientos m\'as avanzados. Las instituciones, tanto p\'ublicas como
privadas, y las asociaciones profesionales tienen la responsabilidad de
abogar por pol\'iticas que promuevan la asequibilidad y la
distribuci\'on justa de estos recursos.

\section*{Conclusi\'on}

Finalmente, es crucial situar la neuromodulaci\'on en su contexto
adecuado. No es una panacea, sino una herramienta m\'as, aunque muy
poderosa, dentro de un modelo de atenci\'on integral y
multidisciplinario. El \'exito terap\'eutico a largo plazo rara vez
depender\'a de una sola modalidad. La combinaci\'on sin\'ergica de la
neuromodulaci\'on con la psicoterapia, la farmacoterapia y las
intervenciones psicosociales probablemente ofrezca los mejores
resultados. El futuro de la psiquiatr\'ia no es una elecci\'on entre
«lo biol\'ogico» y «lo psicosocial», sino su integraci\'on inteligente.

La era de la neuromodulaci\'on ya est\'a aqu\'i. Nuestra tarea como
comunidad profesional es recibirla no con una aceptaci\'on pasiva, sino
con un compromiso activo para desarrollar las competencias, los marcos
\'eticos y los modelos institucionales necesarios para aprovechar su
potencial de manera responsable. Es un desaf\'io formidable, pero
tambi\'en una oportunidad extraordinaria para transformar la vida de
nuestros pacientes y revitalizar nuestra disciplina.

% ── REFERENCIAS ──────────────────────────────────────────────
% Sin \APMrefsbreak. Aqui las referencias arrancan a media columna (a
% 175 pt del tope de la pagina 2, tras el cuerpo) y flushend equilibra la
% ultima pagina por su cuenta: hasta el 1 de septiembre de 2026 habia un
% corte tras la quinta entrada y, medido, no hacia nada (sin el: las
% mismas 154 lineas en las mismas coordenadas y el mismo hash de texto).
% El corte va solo donde caiga el corte de columna, cuando la lista
% arranca al tope de una columna (FM05, LEEME.md).

\APMrefsstart

\APMref{American Psychiatric Association. (2024). APA position statement
  on repetitive transcranial magnetic stimulation (TMS). \textit{APA
  Official Policy.}
  \url{https://www.psychiatry.org/getattachment/23523f4c-6f10-42c6-a4e0-9d54eae5b900/Position-Statement-Transcranial-MagneticStimulation.pdf}}

\APMref{Berlim, M. T., van den Eynde, F., \& Daskalakis, Z. J. (2013).
  Clinically meaningful efficacy and acceptability of low-frequency
  repetitive transcranial magnetic stimulation for treating primary major
  depression: A meta-analysis. \textit{J.~Clin.~Psychiatry, 74}(5),
  e410--e416. \url{https://doi.org/10.4088/JCP.12r07972}}

\APMref{Bikson, M., Brunoni, A. R., Charvet, L. E., et al. (2018).
  Rigor and reproducibility in research with transcranial electrical
  stimulation: An NIMH-supported working group report. \textit{Brain
  Stimulation, 11}(3), 465--480.
  \url{https://doi.org/10.1016/j.brs.2018.01.003}}

\APMref{Fekadu, A., Wooderson, S., Markopoulo, K., et al. (2018).
  Definitions from treatment guidelines and staging methods for
  treatment-resistant depression. \textit{BMC Psychiatry, 18,}
  Art.~100. \url{https://doi.org/10.1186/s12888-018-1679-x}}

\APMref{Fregni, F., El-Hagrassy, M. M., Pacheco-Barrios, K., et al.
  (2021). Evidence-based guidelines and secondary meta-analysis for the
  use of transcranial direct current stimulation in neurological and
  psychiatric disorders. \textit{Int.~J.~Neuropsychopharmacol., 24}(4),
  256--313. \url{https://doi.org/10.1093/ijnp/pyaa051}}

\APMref{Hern\'andez, R. C., Espinosa, M. O., Sampieri-Cabrera, R., \&
  Lara, O. R. A. (2025). Technical report: Efficacy and safety of
  low-intensity transcranial magnetic stimulation in the remission of
  depressive symptoms in patients with treatment-resistant depression in
  Mexico. \textit{Cureus, 17,} Art.~e38832162.
  \url{https://pubmed.ncbi.nlm.nih.gov/38832162/}}

\APMref{Journal of Clinical Psychiatry. (2023). Efficacy and safety of
  transcranial direct current stimulation in major depressive disorder:
  A meta-analysis. \textit{84}(1), 1--10.
  \url{https://pubmed.ncbi.nlm.nih.gov/31837388/}}

\APMref{Lefaucheur, J. P., Aleman, A., Baeken, C., et al. (2020).
  Evidence-based guidelines on the therapeutic use of repetitive
  transcranial magnetic stimulation (rTMS): An update (2014--2018).
  \textit{Clin.~Neurophysiol., 131}(2), 474--528.
  \url{https://doi.org/10.1016/j.clinph.2019.11.002}}

\APMref{McIntyre, R. S., Alsuwaidan, M., Baune, B. T., et al. (2023).
  Treatment-resistant depression: Definition, prevalence, detection,
  management, and investigational interventions. \textit{World
  Psychiatry, 22}(3), 394--412.
  \url{https://doi.org/10.1002/wps.21120}}

\APMref{News-Medical. (2025, March 21). \textit{Nearly half of
  depression patients struggle with treatment resistance.}
  \url{https://www.news-medical.net/news/20250321/}}

\APMref{Padberg, F., \& George, M. S. (2009). Repetitive transcranial
  magnetic stimulation of the prefrontal cortex in depression.
  \textit{Exp.~Neurol., 219}(1), 2--13.
  \url{https://doi.org/10.1016/j.expneurol.2009.04.020}}

\APMref{Ram\'irez-L\'opez, S. A., D\'iaz, V., \& Leticia, S. (2023).
  Extinction therapy and transcranial magnetic stimulation for the
  treatment of anxiety disorders. \textit{Horizonte Sanitario, 22,}
  1--10. \url{https://revistahorizonte.ujat.mx/}}

\APMref{Rush, A. J., Trivedi, M. H., Wisniewski, S. R., et al. (2006).
  Acute and longer-term outcomes in depressed outpatients requiring one
  or several treatment steps: A STAR*D report. \textit{Am.~J.~Psychiatry,
  163}(11), 1905--1917.
  \url{https://doi.org/10.1176/ajp.2006.163.11.1905}}

\APMref{Times Union. (2023, October 10). \textit{This Mexican startup
  makes psychiatric treatment 90\% cheaper, and that's even caught the
  attention of Mark Zuckerberg's sister, who has become his new
  investor.} \url{https://www.timesunion.com/}}

\APMrefsend

% ── COLOPHON ─────────────────────────────────────────────────
{\hfuzz=15pt\APMcolophon}

\end{document}
